% Recuperación documental para el artículo V; RESULTADO_RECUPERADO.
% Propietario: /Users/ruben/Documents/New project/output/REVISION_EDITORIAL_HMT_MD_20260905/01_FUENTES/INTEGRAL/tree/New project/output/ENSAMBLADO_CAUSAL_RESTAURADO_HMT_MD_20260905/fuente/manuscrito/sections/md/delta_127/04b_completaciones_cuanticas.tex
% Líneas de origen: 1--334.
% REV02: construcción algebraica; la aplicación estadística se reúne en 40.
% Correspondencia editorial: gestion/CONSOLIDACION_PAULI_FOCK_REV02.json.
\section{Construction algébrique des complétions quantiques sur APP}
\label{v:rec:chap:completaciones-cuanticas}
\label{v:sec:completaciones-algebraicas}

APP et TPK fournissent des états orientés, un transport et une parité de feuillet.
Leur réalisation linéaire permet de construire les algèbres symétrique et extérieure,
avec leurs opérateurs de création, d’annihilation et d’occupation. Cette section
établit ces constructions algébriques, le secteur cellulaire des champs et
le couplage réversible de mesure. L’application des spectres d’occupation à l’équilibre est développée dans la section~\ref{v:rec:ocupacion:section:02}.
Chaque partie conserve les conditions de la représentation qu’elle utilise.

\subsection{Espace linéaire des états associé à une réalisation de particule}
\label{v:rec:completaciones:section:02}

Soit \(S\) un ensemble fini de classes de chemins TPK à la profondeur choisie
et soit
\[
 \mathcal H_1=\ell^2(S)
\]
l’espace de Hilbert de base orthonormale \((e_s)_{s\in S}\). Un choix de
poids positifs \(E_s\) définit l’opérateur diagonal
\[
 He_s=E_se_s.
\]
Le choix de \(S\), de son quotient par les symétries et des poids fait partie
du modèle. Une fois ce choix fixé, les complétions multiparticulaires sont canoniques.

\subsection{Complétions symétrique et extérieure}
\label{v:rec:completaciones:section:03}

On définit
\[
 \mathcal F_+(\mathcal H_1)=
 \bigoplus_{n\geq0}\operatorname{Sym}^n\mathcal H_1,
 \qquad
 \mathcal F_-(\mathcal H_1)=
 \bigoplus_{n\geq0}\bigwedge^n\mathcal H_1.
\]
La première est la complétion bosonique ; la seconde, la fermionique. Dans les deux,
le terme de somme directe de degré zéro est l’espace du vide. Les puissances extérieures
sont munies du produit scalaire pour lequel les produits extérieurs ordonnés de vecteurs
d’une base orthonormale sont orthonormaux ; les puissances symétriques sont utilisées
avec la normalisation usuelle des opérateurs de création et d’annihilation.

\subsection{Graduation de feuillet et algèbre d’échange}
\label{v:rec:completaciones:section:04}

Supposons que le feuillet orienté définit un caractère additif
\[
 p:\operatorname{Mor}(\mathsf P_{\rm TPK})\longrightarrow\mathbb Z/2\mathbb Z,
 \qquad p(\gamma\eta)=p(\gamma)+p(\eta).
\]
Sur l’espace linéaire gradué correspondant, on introduit la symétrie
\[
 \tau(v\otimes w)=(-1)^{p(v)p(w)}w\otimes v.
\]

\begin{theorem}[Complétion statistique déterminée par la parité]
\label{v:rec:completaciones:statement:02}
Le quotient de l’algèbre tensorielle par les relations
\[
 v\otimes w-(-1)^{p(v)p(w)}w\otimes v
\]
est symétrique sur le secteur pair et extérieur sur le secteur impair. Deux états
impairs identiques ont un produit nul ; les états pairs admettent une occupation
arbitraire.
\end{theorem}

\begin{proof}
Si \(p(v)=p(w)=0\), la relation est \(vw=wv\). Si
\(p(v)=p(w)=1\), elle est \(vw=-wv\) ; en prenant \(v=w\) et en travaillant en
caractéristique zéro, on obtient \(v^2=0\). Les secteurs mixtes commutent sans
signe supplémentaire. C’est la propriété universelle de l’algèbre symétrique
graduée.
\end{proof}

\subsection{Opérateurs canoniques et idempotence de l’occupation}
\label{v:pauli:car-construccion}

Le caractère de parité détermine la complétion ; le produit de cette
complétion détermine ensuite l’algèbre d’opérateurs. Dans le
secteur extérieur, la création \(a_s^\dagger\) est le produit extérieur par
\(e_s\), et \(a_s\) est son adjoint.

\begin{theorem}[Principe d’exclusion de Pauli dans la complétion extérieure]
\label{v:rec:completaciones:statement:01}
\label{v:rec:thm:principio-exclusion-pauli}
Pour tout mode \(s\in S\), les opérateurs de création et d’annihilation de la
complétion extérieure satisfont
\begin{equation}
 \{a_s,a_t^\dagger\}=\delta_{st}I,
 \qquad
 \{a_s,a_t\}=\{a_s^\dagger,a_t^\dagger\}=0.
 \label{v:pauli:eq:car}
\end{equation}
En particulier, l’opérateur d’occupation
\(N_s=a_s^\dagger a_s\) est un projecteur :
\begin{equation}
 N_s^2=N_s,
 \qquad \Spec(N_s)\subseteq\{0,1\}.
 \label{v:pauli:eq:proyeccion-ocupacion}
\end{equation}
Les valeurs propres \(0\) et \(1\) correspondent, respectivement, à l’absence
et à la présence du mode \(s\).
\end{theorem}

\begin{proof}
Sur un produit extérieur ordonné, \(a_s\) supprime le facteur \(e_s\), lorsqu’il est
présent, avec le signe de sa position ; s’il est absent, il donne zéro. Composer
cette opération avec le produit extérieur démontre les relations
\eqref{v:pauli:eq:car} : pour \(s=t\), les deux compositions couvrent les cas
complémentaires d’absence et de présence ; pour \(s\ne t\), les signes
d’échange s’annulent. L’antisymétrie implique
\(e_s\wedge e_s=0\) et
\(a_s^2=(a_s^\dagger)^2=0\). Alors
\begin{equation}
 N_s^2=a_s^\dagger a_s a_s^\dagger a_s
 =a_s^\dagger(I-a_s^\dagger a_s)a_s
 =a_s^\dagger a_s=N_s.
 \label{v:pauli:eq:prueba-idempotencia}
\end{equation}
L’opérateur est autoadjoint, et toute valeur propre d’un idempotent satisfait
\(\lambda^2=\lambda\). Son spectre est donc contenu dans
\(\{0,1\}\). Le vide et l’état \(e_s\) réalisent les deux valeurs
\cite{Pauli1925,Dirac1926}.
\end{proof}

Dans le secteur symétrique, les opérateurs bosoniques \(b_s^\dagger,b_s\)
agissent sur la base normalisée des occupations par les facteurs
\(\sqrt{m_s+1}\) et \(\sqrt{m_s}\), respectivement. De leurs compositions,
on obtient
\begin{equation}
 [b_s,b_t^\dagger]=\delta_{st}I,
 \qquad [b_s,b_t]=[b_s^\dagger,b_t^\dagger]=0.
 \label{v:pauli:eq:ccr}
\end{equation}
Ces identités sont comprises sur le domaine algébrique des particules
finies, stable sous les opérateurs considérés.

Le théorème~\ref{v:rec:completaciones:statement:02} transforme une parité compositionnelle TPK en une règle exacte
d’occupation. Pour obtenir le théorème physique de spin--statistique, il faut
construire en outre une théorie locale lorentzienne positive et démontrer que la
parité de feuillet coïncide avec le spin sous les rotations. L’égalité ne se
déduit pas uniquement du fait que les deux objets prennent leurs valeurs modulo deux.
